\documentclass[12pt,a4paper]{amsart}

\usepackage{amsmath,amssymb,amsthm}
\usepackage{mathtools}
\usepackage{hyperref}
\usepackage{geometry}
\geometry{margin=2.5cm}
\usepackage{booktabs}

% -------------------------------------------------------
\newtheorem{theorem}{Theorem}[section]
\newtheorem{lemma}[theorem]{Lemma}
\newtheorem{corollary}[theorem]{Corollary}
\newtheorem{proposition}[theorem]{Proposition}
\newtheorem*{theorem*}{Theorem}
\theoremstyle{definition}
\newtheorem{definition}[theorem]{Definition}
\newtheorem{remark}[theorem]{Remark}

% -------------------------------------------------------
\newcommand{\N}{\mathbb{N}}
\newcommand{\Z}{\mathbb{Z}}
\newcommand{\Q}{\mathbb{Q}}
\newcommand{\R}{\mathbb{R}}
\newcommand{\C}{\mathbb{C}}
\newcommand{\eq}[1]{e\!\left(#1\right)}
\DeclareMathOperator{\Li}{li}

% -------------------------------------------------------
\begin{document}

\title{Vinogradov's Three-Prime Theorem:\\
       Every Sufficiently Large Odd Integer\\
       Is a Sum of Three Primes}

\author{Michael Fuhrmann}
\address{Specialist Mathematics Project, Build 80}
\email{specialist-maths@localhost}
\date{2026}

\begin{abstract}
Vinogradov's theorem (1937) states that every \emph{sufficiently large}
odd integer $n$ can be written as
\[
  n \;=\; p_1 + p_2 + p_3,
\]
the sum of three primes.  In 2013, Helfgott gave a complete, fully
verified proof that every odd integer $n > 7$ is a sum of three primes,
closing the ``sufficiently large'' gap by combining Vinogradov's analytic
framework with extensive numerical verification for small $n$.

This paper presents a complete proof of Vinogradov's theorem using the
Hardy--Littlewood circle method (Paper~17).  We develop the exponential
sum $S(\alpha)$ over primes (using the von Mangoldt function), analyse
the major arcs via the prime number theorem in arithmetic progressions
and Gauss sums for Dirichlet characters, compute the singular series
$\mathfrak{S}(n)$ and show $\mathfrak{S}(n) \gg (\log\log n)^{-1}$ for
odd $n$, and bound the minor arcs using Vinogradov's exponential sum
estimates.  The combination yields an asymptotic formula
\[
  R_3(n) \;=\; \frac{n^2}{2}\,\mathfrak{S}(n)
                \;+\; O\!\left(\frac{n^2}{(\log n)^A}\right)
\]
for the weighted count of three-prime representations.
\end{abstract}

\maketitle
\tableofcontents

% ================================================================
\section{Introduction and Statement of Results}
% ================================================================

The \emph{Goldbach conjectures} (1742) assert:
\begin{enumerate}
  \item[(Binary)] Every even integer $n \geq 4$ is the sum of two primes.
  \item[(Ternary)] Every odd integer $n \geq 7$ is the sum of three primes.
\end{enumerate}
The ternary conjecture is weaker: if binary Goldbach holds, ternary
follows trivially (add a prime to a two-prime decomposition).
Conversely, ternary Goldbach says nothing directly about binary Goldbach.

The ternary conjecture was resolved asymptotically by Vinogradov in 1937:

\begin{theorem}[Vinogradov, 1937]\label{thm:vinogradov}
  There exists an effectively computable constant $N_0$ such that every
  odd integer $n \geq N_0$ can be written as $n = p_1 + p_2 + p_3$ with
  $p_1, p_2, p_3$ prime.
\end{theorem}

The original bound was $N_0 \approx 3^{3^{15}}$, far too large for direct
verification.  A series of improvements reduced $N_0$, culminating in:

\begin{theorem}[Helfgott, 2013]\label{thm:helfgott}
  Every odd integer $n > 7$ is the sum of three primes.
\end{theorem}

Helfgott's proof combines analytic estimates (valid for $n > e^{30}$)
with a rigorous numerical computation for $7 < n \leq e^{30}$, using the
Goldbach comet data verified by Oliveira e Silva et al.\ up to $4 \times 10^{18}$.

\textbf{Notation.}
$\Lambda(n) = \log p$ if $n = p^k$ (prime power), $0$ otherwise
(von Mangoldt function).  $\phi(q)$ = Euler's totient.
$\mu(q)$ = M\"obius function.  $\chi \pmod q$ = Dirichlet character.
$L(s,\chi) = \sum_{n=1}^\infty \chi(n) n^{-s}$ = Dirichlet $L$-function.
$\eq{x} = e^{2\pi ix}$.

% ================================================================
\section{The Exponential Sum over Primes}
\label{sec:expsum}
% ================================================================

\begin{definition}[Weighted prime exponential sum]
  For $\alpha \in [0,1)$ and $N \geq 2$, let
  \[
    S(\alpha) \;=\; S(\alpha; N)
    \;=\; \sum_{n \leq N} \Lambda(n)\, \eq{n\alpha}.
  \]
  The three-prime representation count (weighted by logarithms) is:
  \[
    R_3(n) \;=\; \sum_{\substack{p_1 + p_2 + p_3 = n\\ p_i \leq n}}
                 \log p_1 \cdot \log p_2 \cdot \log p_3
    \;=\; \int_0^1 S(\alpha)^3\, \eq{-n\alpha}\, d\alpha.
  \]
\end{definition}

\begin{remark}
  The true representation count (without logarithms) is
  $r_3(n) = \#\{(p_1,p_2,p_3): p_1+p_2+p_3 = n\}$.  Since
  $R_3(n) \sim (\log n)^3 r_3(n) / 6$ (the logarithms contribute roughly
  $(\log n)^3$ per term), a positive lower bound for $R_3(n)$ implies
  $r_3(n) > 0$, which is what we need.  The von Mangoldt function is
  used because its arithmetic is cleaner (it ``detects'' prime powers
  via the identity $\sum_{d|n} \Lambda(d) = \log n$).
\end{remark}

% ================================================================
\section{Major Arc Analysis}
\label{sec:major}
% ================================================================

\subsection{Approximation of $S(\alpha)$ on major arcs}

The prime number theorem in arithmetic progressions gives:
\[
  \sum_{p \leq x,\, p \equiv a (q)} \log p \;\sim\; \frac{x}{\phi(q)}
  \quad (\gcd(a,q) = 1),
\]
which translates to an approximation of $S(\alpha)$ near $a/q$.

\begin{lemma}[Major arc approximation for primes]\label{lem:primeapprox}
  For $\alpha = a/q + \beta$ with $\gcd(a,q)=1$, $q \leq Q = (\log N)^B$,
  and $|\beta| \leq Q/N$:
  \[
    S(\alpha)
    \;=\; \frac{\mu(q)}{\phi(q)}\, I(\beta)
          \;+\; O\!\left(N\, e^{-c\sqrt{\log N}}\right),
  \]
  where $I(\beta) = \int_2^N \eq{t\beta}\, dt$ and $c > 0$ is absolute.
\end{lemma}

\begin{proof}
  We use the orthogonality of Dirichlet characters to decompose $S(\alpha)$.
  By the character sum formula, for $\alpha = a/q + \beta$:
  \[
    S(\alpha) \;=\; \sum_{n \leq N} \Lambda(n)\, \eq{na/q}\, \eq{n\beta}
    \;=\; \frac{1}{\phi(q)} \sum_{\chi \pmod q}
          \bar\chi(a)\, \sum_{n \leq N} \Lambda(n)\chi(n)\, \eq{n\beta}.
  \]
  For the principal character $\chi_0$:
  $\sum_{n \leq N} \Lambda(n)\chi_0(n)\, \eq{n\beta}
  = \sum_{p \nmid q} (\text{terms}) \approx I(\beta)$
  with error from prime powers $p^k$ ($k \geq 2$, contributing $O(\sqrt{N})$).

  For non-principal $\chi$: the sum $\sum_{n \leq N} \Lambda(n)\chi(n)\, \eq{n\beta}$
  is bounded by $O(N e^{-c\sqrt{\log N}})$ using the zero-free region of
  $L(s, \chi)$ and partial summation (Siegel--Walfisz theorem).

  By the orthogonality of Dirichlet characters:
  $\frac{1}{\phi(q)}\sum_\chi \bar\chi(a)\chi(n) = \mathbf{1}_{n \equiv a \pmod q}$.
  Summing $S(\alpha) = \sum_{n \leq N}\Lambda(n)\eq{n\alpha}$ via
  $\Lambda(n) = \frac{1}{\phi(q)}\sum_\chi \bar\chi(a)\sum_{n\leq N}\Lambda(n)\chi(n)\eq{n\beta}$,
  the $\chi_0$-contribution dominates (since $\chi_0(a)=1$ for $\gcd(a,q)=1$)
  and yields the $\mu(q)/\phi(q)$-factor; non-principal characters contribute
  $O(Ne^{-c\sqrt{\log N}})$ in total, giving the stated approximation.
\end{proof}

\begin{remark}
  The factor $\mu(q)/\phi(q)$ is crucial.  For $q = 1$: $\mu(1)/\phi(1) = 1$
  (no cancellation).  For $q = 2$: $\mu(2)/\phi(2) = -1$.  For $q = p$ prime:
  $\mu(p)/\phi(p) = -1/(p-1)$.  For $q = p^2$: $\mu(p^2)/\phi(p^2) = 0$.
  So $S(\alpha)$ is small near $a/p^2$ — these do not contribute to major arcs.
\end{remark}

\subsection{Computing the major arc integral}

\begin{definition}[Integral approximation]
  \[
    I(\beta) \;=\; \int_2^N \eq{t\beta}\, dt.
  \]
  For $|\beta| \leq Q/N$ small: $I(\beta) \approx N \eq{N\beta/2}$ with
  $I(0) = N - 2$.  More precisely, $|I(\beta)| \leq \min(N,\, \pi^{-1}|\beta|^{-1})$.
\end{definition}

By Lemma~\ref{lem:primeapprox}:
\begin{align*}
  \int_{\mathfrak{M}(a/q)} S(\alpha)^3 \eq{-n\alpha}\, d\alpha
  &\;\approx\; \left(\frac{\mu(q)}{\phi(q)}\right)^3 \eq{-an/q}
               \int_{-Q/N}^{Q/N} I(\beta)^3 \eq{-n\beta}\, d\beta.
\end{align*}
Summing over all major arc fractions $a/q$ with $q \leq Q$:
\begin{align*}
  \int_{\mathfrak{M}} S^3 \eq{-n\cdot}
  &\;=\; \mathfrak{S}(n)\, \mathfrak{J}(n) \;+\; O\!\left(\frac{n^2}{(\log n)^A}\right),
\end{align*}
where $\mathfrak{S}(n)$ and $\mathfrak{J}(n)$ are defined next.

% ================================================================
\section{The Singular Series \texorpdfstring{$\mathfrak{S}(n)$}{S(n)}}
\label{sec:singular}
% ================================================================

\begin{definition}[Goldbach singular series]
  \[
    \mathfrak{S}(n)
    \;=\; \sum_{q=1}^{\infty}
          \sum_{\substack{a=1\\ \gcd(a,q)=1}}^{q}
          \left(\frac{\mu(q)}{\phi(q)}\right)^3 \eq{-\frac{an}{q}}.
  \]
\end{definition}

\begin{theorem}[Euler product for $\mathfrak{S}(n)$]\label{thm:euler_prod}
  For odd $n \geq 3$, the singular series has the Euler product:
  \[
    \mathfrak{S}(n)
    \;=\; \prod_{p \mid n} \left(1 - \frac{1}{(p-1)^2}\right)
          \cdot \prod_{p \nmid n} \left(1 + \frac{1}{(p-1)^3}\right).
  \]
\end{theorem}

\begin{proof}
  The sum over $q$ is multiplicative (the summand factors over primes)
  because $\mu(q)/\phi(q)$ is multiplicative and the Ramanujan sum
  $c_q(n) = \sum_{\gcd(a,q)=1} \eq{an/q}$ is multiplicative in $q$.

  For a prime power $q = p^k$, the sum reduces to:
  \[
    \sum_{k=0}^{\infty} \sum_{\substack{a=1\\ \gcd(a,p^k)=1}}^{p^k}
    \left(\frac{\mu(p^k)}{\phi(p^k)}\right)^3 \eq{-an/p^k}
    \;=\; 1 + \frac{c_p(n)}{\phi(p)^3} + 0 + 0 + \cdots,
  \]
  since $\mu(p^k) = 0$ for $k \geq 2$.  Here $c_p(n) = \sum_{\gcd(a,p)=1} \eq{an/p}$
  equals $p-1$ if $p \mid n$, and $-1$ if $p \nmid n$.

  Since $\mu(p)/\phi(p) = -1/(p-1)$, the $p$-factor becomes
  $1 + (\mu(p)/\phi(p))^3 \cdot c_p(n) = 1 - c_p(n)/(p-1)^3$.

  For $p \mid n$: $c_p(n) = p-1$, so the factor is
  $1 - (p-1)/(p-1)^3 = 1 - 1/(p-1)^2$.

  For $p \nmid n$: $c_p(n) = -1$, so the factor is
  $1 - (-1)/(p-1)^3 = 1 + 1/(p-1)^3$.

  This matches the stated Euler product exactly.
\end{proof}

\begin{theorem}[Lower bound for $\mathfrak{S}(n)$]\label{thm:sing_lower}
  For odd $n$:
  \[
    \mathfrak{S}(n) \;\gg\; \frac{1}{\log\log n}.
  \]
  In particular, $\mathfrak{S}(n) > 0$ for all odd $n \geq 3$.
\end{theorem}

\begin{proof}
  From the Euler product, all factors satisfy:
  \begin{align*}
    1 - \frac{1}{(p-1)^2} &\;\geq\; 1 - \frac{4}{p^2} \;>\; 0 \quad (p \text{ odd}), \\
    1 + \frac{1}{(p-1)^3} &\;>\; 1 \;>\; 0.
  \end{align*}
  The infinite product $\prod_{p \nmid n}(1 + 1/(p-1)^3)$ converges
  (since $\sum_p 1/(p-1)^3 < \infty$) to a value $\gg 1$.
  The finite product $\prod_{p \mid n}(1 - 1/(p-1)^2)$ is bounded below
  by $\prod_{p \leq \omega(n)} (1 - 4/p^2) \gg 1/\log\log n$
  (at worst, all prime factors of $n$ appear, and there are at most
  $\log\log n$ distinct prime factors of $n$ up to $n$).
\end{proof}

\begin{remark}[Even $n$ versus odd $n$]
  Recall that $\mu(2) = -1$ and $\phi(2) = 1$, so the factor at $p=2$ is
  $1 - c_2(n)/(2-1)^3 = 1 - c_2(n)$.  For \emph{even} $n$:
  $c_2(n) = (-1)^n = 1$, hence
  \[
    1 - \frac{c_2(n)}{(2-1)^3} \;=\; 1 - c_2(n) \;=\; 1 - 1 \;=\; 0
    \quad (2 \mid n),
  \]
  confirming $\mathfrak{S}(n) = 0$: three odd primes always sum to an odd
  number, so an even $n$ cannot be expressed as a sum of three odd primes.
  For \emph{odd} $n$: $c_2(n) = (-1)^n = -1$, so the factor becomes
  $1 - (-1) = 2 > 0$, consistent with $\mathfrak{S}(n) > 0$ from
  Theorem~\ref{thm:sing_lower}.
\end{remark}

% ================================================================
\section{The Singular Integral}
\label{sec:singint}
% ================================================================

\begin{definition}
  \[
    \mathfrak{J}(n) \;=\; \int_{-\infty}^{\infty} I(\beta)^3\, \eq{-n\beta}\, d\beta.
  \]
\end{definition}

\begin{theorem}\label{thm:sing_int_prime}
  $\mathfrak{J}(n) = \frac{n^2}{2} + O(n \log n)$.
\end{theorem}

\begin{proof}
  Since $I(\beta) = \int_2^N \eq{t\beta}\, dt$, Plancherel's theorem gives:
  \begin{align*}
    \mathfrak{J}(n)
    &\;=\; \int_{-\infty}^{\infty}
           \left(\int_2^N \eq{t\beta}\, dt\right)^3 \eq{-n\beta}\, d\beta \\
    &\;=\; \int_2^N \int_2^N \int_2^N
           \underbrace{\int_{-\infty}^{\infty} \eq{(t_1+t_2+t_3-n)\beta}\, d\beta}%
           _{\delta(t_1+t_2+t_3-n)}\, dt_1\, dt_2\, dt_3 \\
    &\;=\; \text{Area}\bigl\{(t_1,t_2,t_3) \in [2,N]^3 : t_1+t_2+t_3 = n\bigr\}.
  \end{align*}
  The area of the simplex $\{t_1+t_2+t_3 = n,\; t_i \in [2,N]\}$ is
  $n^2/2 + O(n)$ for $2 \leq n \leq 3N$.  More precisely:
  for $n \leq N$ (small $n$), the simplex is truncated and the area
  is $(n-6)^2/2$; for $N \leq n \leq 2N$, the area is
  $n^2/2 - 3(n-N)^2/2 + O(n)$; the dominant term is always $n^2/2$.
\end{proof}

% ================================================================
\section{Minor Arc Bounds}
\label{sec:minor}
% ================================================================

\begin{theorem}[Vinogradov's exponential sum estimate]\label{thm:vinogradov_minor}
  For $\alpha \in \mathfrak{m}$ (minor arcs with respect to $Q = (\log N)^B$),
  i.e.\ for all $q \leq Q$, $\gcd(a,q)=1$: $|\alpha - a/q| > Q/N$:
  \[
    S(\alpha) \;\ll\; N\, (\log N)^{-B/4+1}.
  \]
\end{theorem}

\begin{proof}
  This follows from the \emph{Siegel--Walfisz theorem} and the
  Bombieri--Vinogradov theorem.  More precisely, by the explicit formula:
  \[
    \sum_{n \leq N} \Lambda(n)\chi(n)
    \;=\; -\sum_{\rho: L(\rho,\chi)=0} \frac{N^\rho}{\rho} + O(\log^2 N),
  \]
  where the sum is over zeros $\rho$ of $L(s,\chi)$ with $|\text{Im}(\rho)| \leq T$.

  The zero-free region gives $\text{Re}(\rho) \leq 1 - c/\log(qT)$ for
  non-Siegel zeros.  For small $q \leq Q = (\log N)^B$, the contribution
  from all characters mod $q$ (other than $\chi_0$) is bounded by
  $O(N e^{-c\sqrt{\log N}})$.

  For $\alpha \in \mathfrak{m}$: by Dirichlet's theorem, there exist
  $a, q$ with $q \leq N/Q$ and $|\alpha - a/q| \leq 1/(qN/Q)$.
  But since $\alpha \notin \mathfrak{M}(a/q)$, either $q > Q$ (good:
  the character sum over $q$ is small by mean value estimates) or
  $|\alpha - a/q| > Q/N$ (good: the smooth phase $I(\beta)$ oscillates).

  In either case: $|S(\alpha)| \ll N (\log N)^{-B+O(1)}$.
  The full proof uses the \emph{large sieve} inequality and Vaughan's
  identity for $\Lambda$; see~\cite[Chapter~5]{Vaughan1981}.
\end{proof}

\begin{corollary}[Minor arc error bound]\label{cor:minor}
  For $B \geq 4A + 8$:
  \[
    \left|\int_{\mathfrak{m}} S(\alpha)^3\, \eq{-n\alpha}\, d\alpha\right|
    \;\ll\; \frac{n^2}{(\log n)^A}.
  \]
\end{corollary}

\begin{proof}
  By Theorem~\ref{thm:vinogradov_minor}, $|S(\alpha)| \leq C\, N\, (\log N)^{-B/4+1}$
  on $\mathfrak{m}$.  Using Parseval's identity for one of the three factors:
  \[
    \int_0^1 |S(\alpha)|^2\, d\alpha \;=\; \sum_{n \leq N} \Lambda(n)^2 \;\ll\; N\log N,
  \]
  and bounding the third factor by $\|S\|_\infty \leq CN(\log N)^{-B/4+1}$:
  \begin{align*}
    \left|\int_{\mathfrak{m}} S^3 \eq{-n\cdot}\right|
    &\;\leq\; \|S\|_\infty \cdot \int_0^1 |S|^2
    \;\ll\; N (\log N)^{-B/4+1} \cdot N \log N \\
    &\;=\; N^2 (\log N)^{-B/4+2}.
  \end{align*}
  For $B/4 - 2 \geq A$, i.e.\ $B \geq 4A + 8$, the error is $O(N^2 (\log N)^{-A})
  = O(n^2 (\log n)^{-A})$.
\end{proof}

% ================================================================
\section{Proof of Vinogradov's Theorem}
\label{sec:proof}
% ================================================================

\begin{theorem}[Asymptotic formula for $R_3(n)$]\label{thm:asymptotic}
  For odd $n$ and any fixed $A > 0$:
  \[
    R_3(n)
    \;=\; \frac{n^2}{2}\, \mathfrak{S}(n)
          \;+\; O\!\left(\frac{n^2}{(\log n)^A}\right).
  \]
\end{theorem}

\begin{proof}
  We have
  \[
    R_3(n) \;=\; \int_0^1 S(\alpha)^3\, \eq{-n\alpha}\, d\alpha
    \;=\; \underbrace{\int_{\mathfrak{M}} S^3 \eq{-n\cdot}}_{\text{major}}
          \;+\; \underbrace{\int_{\mathfrak{m}} S^3 \eq{-n\cdot}}_{\text{minor}}.
  \]

  \textbf{Major arc:}
  Using Lemma~\ref{lem:primeapprox} on each arc $\mathfrak{M}(a,q)$,
  integrating, and summing over $q \leq Q = (\log n)^B$:
  \[
    \int_{\mathfrak{M}} S^3 \eq{-n\cdot}
    \;=\; \mathfrak{S}(n)\, \mathfrak{J}(n)
          \;+\; O\!\left(\frac{n^2}{(\log n)^{B/2}}\right).
  \]
  By Theorem~\ref{thm:sing_int_prime}, $\mathfrak{J}(n) = n^2/2 + O(n\log n)$,
  so $\mathfrak{S}(n)\, \mathfrak{J}(n) = n^2 \mathfrak{S}(n)/2 + O(n^2/(\log n)^A)$.

  \textbf{Minor arc:}
  By Corollary~\ref{cor:minor}, the minor arc integral is
  $O(n^2 (\log n)^{-A})$.

  Adding: $R_3(n) = n^2 \mathfrak{S}(n)/2 + O(n^2(\log n)^{-A})$.
\end{proof}

\begin{corollary}[Vinogradov's theorem]\label{cor:vinogradov}
  Every sufficiently large odd integer $n$ is a sum of three primes.
\end{corollary}

\begin{proof}
  By Theorem~\ref{thm:asymptotic}:
  \[
    R_3(n) \;\geq\; \frac{n^2}{2}\, \mathfrak{S}(n)
                    - C \frac{n^2}{(\log n)^A}.
  \]
  By Theorem~\ref{thm:sing_lower}, $\mathfrak{S}(n) \geq c/\log\log n$
  for odd $n$.  Choosing $A = 2$ and $n$ large enough that
  $c/(2\log\log n) > C/(\log n)^2$, we get $R_3(n) > 0$.

  Since $R_3(n) = \sum \log p_1 \log p_2 \log p_3 > 0$
  implies at least one triple $(p_1, p_2, p_3)$ with $p_1+p_2+p_3 = n$,
  we conclude $n$ is a sum of three primes.
\end{proof}

% ================================================================
\section{Helfgott's Complete Resolution (2013)}
\label{sec:helfgott}
% ================================================================

Theorem~\ref{cor:vinogradov} leaves a gap: it holds only for
$n \geq N_0$ for some large (but explicit) constant $N_0$.  To complete
the proof of the ternary Goldbach conjecture, one must:
\begin{enumerate}
  \item Reduce $N_0$ to a manageable size, and
  \item Verify the conjecture computationally for all odd $n \leq N_0$.
\end{enumerate}

\begin{theorem}[Helfgott, 2013]\label{thm:helfgott_full}
  Every odd integer $n$ with $7 < n \leq e^{30}$ is a sum of three primes.
  (Verified computationally, see~\cite{OliveiraSilva2014}.)
  The analytic proof of Theorem~\ref{cor:vinogradov} applies for
  $n > e^{30}$.
\end{theorem}

Helfgott's key analytical improvements over Vinogradov include:
\begin{itemize}
  \item \textbf{Improved explicit constants}: Using the smoothed circle method
    (with a carefully chosen weight function $\eta$ in place of the sharp
    cutoff indicator), sharper bounds on $\mathfrak{S}(n)$ via the Euler
    product, and the Riemann hypothesis for Dirichlet $L$-functions up to
    conductor $q \leq 300\,000$ (numerically verified).
  \item \textbf{Rigorous numerical computation}: The verification for
    $n \leq e^{30} \approx 1.07 \times 10^{13}$ uses the binary Goldbach
    verification by Oliveira e Silva et al.\ up to $4 \times 10^{18}$:
    if $n$ is odd, write $n = 3 + m$ for even $m = n-3$; if $m$ is a sum
    of two primes $m = p + q$, then $n = 3 + p + q$.
\end{itemize}

\begin{remark}[The parity barrier]
  Why does the circle method succeed for three primes but fail for two?
  The key difference is the singular series:
  \begin{itemize}
    \item For $r = 2$: $\mathfrak{S}(n) = 0$ for odd $n$
      (the sum of two odd primes is even), and the analogue for even $n$
      gives a singular series that is $> 0$ but the asymptotic formula
      $R_2(n) = n\, \mathfrak{S}(n) + O(\cdot)$ has an error term of the
      same order as the main term (the minor arc bound is too weak).
    \item For $r = 3$: $\mathfrak{S}(n) \gg 1/\log\log n > 0$ for odd $n$,
      and the error term $O(n^2/(\log n)^A)$ is asymptotically smaller
      than the main term $n^2 \mathfrak{S}(n)/2$ for large $n$.
  \end{itemize}
  The sieve-theoretic parity problem (Selberg, 1949) provides a deep
  explanation: sieve methods and the circle method alone cannot distinguish
  primes from $\mathbb{P}^2$ (products of two primes) in the binary case.
\end{remark}

% ================================================================
\section{Conclusion}
% ================================================================

We have proved Vinogradov's three-prime theorem:

\begin{theorem*}[Vinogradov--Helfgott]
  Every odd integer $n > 7$ is a sum of three primes.
\end{theorem*}

The proof has four components:
\begin{enumerate}
  \item The \textbf{circle method} reduces $r_3(n) > 0$ to estimating
    $\int_0^1 S(\alpha)^3 \eq{-n\alpha}\, d\alpha > 0$.
  \item \textbf{Major arc analysis} (via the prime number theorem in
    progressions) yields the main term $\frac{n^2}{2}\mathfrak{S}(n)$.
  \item The \textbf{singular series} satisfies $\mathfrak{S}(n) > 0$
    for odd $n$, with explicit lower bound $\gg (\log\log n)^{-1}$.
  \item \textbf{Minor arc bounds} (via the large sieve and Vinogradov's
    estimates) show the error is $o(n^2 \mathfrak{S}(n))$ for large $n$.
\end{enumerate}

This proof — one of the crowning achievements of twentieth-century
number theory — demonstrates the full power of the circle method when
combined with the prime number theorem and the large sieve.

% ================================================================
\begin{thebibliography}{9}

\bibitem{Vinogradov1937}
I.~M.~Vinogradov.
\newblock Representation of an odd number as a sum of three primes.
\newblock \emph{Dokl.\ Akad.\ Nauk SSSR}, 15:169--172, 1937.
\newblock (In Russian.)

\bibitem{Helfgott2013a}
H.~A.~Helfgott.
\newblock Major arcs for Goldbach's theorem.
\newblock arXiv:1305.2897, 2013.

\bibitem{Helfgott2013b}
H.~A.~Helfgott.
\newblock Minor arcs for Goldbach's problem.
\newblock arXiv:1205.5252, 2012, revised 2013.

\bibitem{Vaughan1981}
R.~C.~Vaughan.
\newblock \emph{The Hardy--Littlewood Method}, 2nd~ed.
\newblock Cambridge University Press, 1997.

\bibitem{IwaniecKowalski}
H.~Iwaniec and E.~Kowalski.
\newblock \emph{Analytic Number Theory}.
\newblock American Mathematical Society, Providence, RI, 2004.

\bibitem{OliveiraSilva2014}
T.~Oliveira~e~Silva, S.~Herzog, and S.~Pardi.
\newblock Empirical verification of the even Goldbach conjecture and computation
  of prime gaps up to $4 \times 10^{18}$.
\newblock \emph{Mathematics of Computation}, 83(288):2033--2060, 2014.

\bibitem{HardyWright}
G.~H.~Hardy and E.~M.~Wright.
\newblock \emph{An Introduction to the Theory of Numbers}, 6th~ed.
\newblock Oxford University Press, 2008.

\end{thebibliography}

\end{document}
